% généré depuis source/08_conclusion_generale.md — ne pas modifier à la main
\chapter*{Conclusion générale}
\addcontentsline{toc}{chapter}{Conclusion générale}
\frontmark{Conclusion générale}

Ce projet a appliqué l'ensemble de la démarche de sûreté de fonctionnement à une flotte de turboréacteurs simulés, des données brutes jusqu'à un outil d'aide à la décision.

L'analyse de fiabilité, menée en tenant compte des 100 moteurs censurés, montre une flotte soumise à l'usure, avec une période sans défaillance. La loi de Weibull à trois paramètres, acceptée par le test KS et équivalente à la log-normale selon l'AIC, donne un MTTF de 206,3 cycles et des durées $B_{10}$ et $B_{50}$ de 151,9 et 199,9 cycles. Sous des hypothèses de réparation et de coût explicitement signalées, la disponibilité intrinsèque atteint 95,98 \% et un remplacement préventif à 131 cycles réduit le coût par cycle de 68 \%.

La chaîne de maintenance prédictive complète cette vision par une décision moteur par moteur. Un HI construit sur 12 capteurs décrit la dégradation, le classifieur détecte l'état critique avec une AUC de 0,991, et deux modèles estiment la RUL. Le LSTM est le plus précis (RMSE de 15,35 cycles) ; le modèle de dégradation bayésien, moins précis (RMSE de 24,47 cycles), est interprétable et probabiliste, et son biais comme la calibration insuffisante de ses intervalles sont mesurés et expliqués.

Les résultats sont restitués dans un tableau de bord et produits par un système reproductible. Les travaux futurs prioritaires portent sur un modèle de dégradation plus fidèle et mieux calibré, l'ajustement de la loi retenue avec censure et l'extension à plusieurs conditions opératoires.
